% GENERATED FILE. Edit canonical lessons, then run npm run generate:books.
% canonical-source-hash: fnv1a64:65a8fe1f8045523f
% canonical-lessons: HI-C333-ghat, HI-C333-visarjan, HI-C333-dubona, HI-C333-bikharna, HI-C333-dipdan

\chapter{Down to the River}
\label{ch:hi-down-to-the-river}
\hlchaptermodality{333}

\begin{chapteropening}
\textbf{By the end of this chapter:} \emph{I can say what people did at the river on a festival.}

Complete the last lesson of chapter 333: I can say what people did at the river on a festival.
\end{chapteropening}

\section[\texorpdfstring{\hi{घाट} (ghāṭ)}{घाट (ghāṭ)}]{\hi{घाट} (ghāṭ) — a ghat, steps down to a river}

\label{lesson:HI-C333-ghat}



\begin{quote}
Before the new one: say the Hindi for a bonfire, then the Hindi for an effigy.
\end{quote}

\subsection*{You'll want to know: \hi{घाट}}
\textbf{\hi{घाट}} — \emph{ghāṭ} — \textquotedblleft{}a ghat, steps down to a river\textquotedblright{}.

\begin{quote}
\hi{माफ़} \hi{कीजिए}, \hi{घाट} \hi{कैसे} \hi{जाते} \hi{हैं}?
\end{quote}

\begin{quote}
\emph{māf kījiye, ghāṭ kaise jāte haiṁ?}
\end{quote}

\begin{quote}
Excuse me, how do I get to the ghat?
\end{quote}

Masculine.

\begin{cousinweb}
\textbf{\hi{घाट}} (\emph{ghāṭ}) goes back to Sanskrit \emph{ghaṭṭa}, \textquotedblleft{}a landing place\textquotedblright{} (CDIAL 4414): the steps where people come down to bathe.
\end{cousinweb}

\begin{grammarlens}[title={What you've built}]
One more word for telling what happened.
\end{grammarlens}

\subsection*{Your turn}
\begin{itemize}
\raggedright
  \item \emph{Say it:} \emph{ghāṭ}
  \item \emph{Say it:} \emph{ghāṭ}, once more
  \item \emph{Say it:} \emph{māf kījiye, ghāṭ kaise jāte haiṁ?} — Excuse me, how do I get to the ghat?
  \item \emph{Recall:} say the Hindi for a bonfire, then the Hindi for an effigy, then say \emph{ghāṭ} again
\end{itemize}

\begin{tcolorbox}[breakable,colback=teal!4,colframe=teal!35!black,title={Before you move on}]
What does \hi{घाट} mean? (\textquotedblleft{}A ghat, steps down to a river\textquotedblright{}.) Say it once more.
\end{tcolorbox}
\section[\texorpdfstring{\hi{विसर्जन} (visarjan)}{विसर्जन (visarjan)}]{\hi{विसर्जन} (visarjan) — immersion (of an image)}

\label{lesson:HI-C333-visarjan}



\begin{quote}
Before the new one: say the Hindi for an effigy, then the Hindi for a ghat, steps down to a river.
\end{quote}

\subsection*{You'll want to know: \hi{विसर्जन}}
\textbf{\hi{विसर्जन}} — \emph{visarjan} — \textquotedblleft{}immersion (of an image)\textquotedblright{}.

\begin{quote}
\hi{शाम} \hi{को} \hi{नदी} \hi{में} \hi{मूर्ति} \hi{का} \hi{विसर्जन} \hi{हुआ।}
\end{quote}

\begin{quote}
\emph{śām ko nadī meṁ mūrti kā visarjan huā.}
\end{quote}

\begin{quote}
In the evening the image was immersed in the river.
\end{quote}

Masculine.

\begin{cousinweb}
\textbf{\hi{विसर्जन}} (\emph{visarjan}) is Sanskrit \emph{visarjana}, \textquotedblleft{}a letting go, a sending away\textquotedblright{} (Monier-Williams): at the end of a festival the image is let go into the water.
\end{cousinweb}

\begin{grammarlens}[title={What you've built}]
One more word for telling what happened.
\end{grammarlens}

\subsection*{Your turn}
\begin{itemize}
\raggedright
  \item \emph{Say it:} \emph{visarjan}
  \item \emph{Say it:} \emph{visarjan}, once more
  \item \emph{Say it:} \emph{śām ko nadī meṁ mūrti kā visarjan huā.} — In the evening the image was immersed in the river.
  \item \emph{Recall:} say the Hindi for an effigy, then the Hindi for a ghat, steps down to a river, then say \emph{visarjan} again
\end{itemize}

\begin{tcolorbox}[breakable,colback=teal!4,colframe=teal!35!black,title={Before you move on}]
What does \hi{विसर्जन} mean? (\textquotedblleft{}Immersion (of an image)\textquotedblright{}.) Say it once more.
\end{tcolorbox}
\section[\texorpdfstring{\hi{डुबोना} (ḍubonā)}{डुबोना (ḍubonā)}]{\hi{डुबोना} (ḍubonā) — to dip; to immerse}

\label{lesson:HI-C333-dubona}



\begin{quote}
Before the new one: say the Hindi for a ghat, steps down to a river, then the Hindi for immersion (of an image).
\end{quote}

\subsection*{You'll want to know: \hi{डुबोना}}
\textbf{\hi{डुबोना}} — \emph{ḍubonā} — \textquotedblleft{}to dip; to immerse\textquotedblright{}.

\begin{quote}
\hi{भक्तों} \hi{ने} \hi{फूल} \hi{पानी} \hi{में} \hi{डुबोए।}
\end{quote}

\begin{quote}
\emph{bhaktoṁ ne phūl pānī meṁ ḍuboe.}
\end{quote}

\begin{quote}
The devotees dipped flowers in the water.
\end{quote}

Past, with \hi{ने} (he or she …): \hi{उसने} \hi{कपड़ा} \hi{पानी} \hi{में} \hi{डुबोया}, \emph{usne kapṛā pānī meṁ ḍuboyā}.

\begin{cousinweb}
\textbf{\hi{डुबोना}} (\emph{ḍubonā}) is the causative of \textbf{\hi{डूबना}}, to sink, which CDIAL 5561 sets under \emph{*ḍubb}: to make something go under the water.
\end{cousinweb}

\begin{grammarlens}[title={What you've built}]
One more word for telling what happened.
\end{grammarlens}

\subsection*{Your turn}
\begin{itemize}
\raggedright
  \item \emph{Say it:} \emph{ḍubonā}
  \item \emph{Say it:} \emph{ḍubonā}, once more
  \item \emph{Say it:} \emph{bhaktoṁ ne phūl pānī meṁ ḍuboe.} — The devotees dipped flowers in the water.
  \item \emph{Recall:} say the Hindi for a ghat, steps down to a river, then the Hindi for immersion (of an image), then say \emph{ḍubonā} again
\end{itemize}

\begin{tcolorbox}[breakable,colback=teal!4,colframe=teal!35!black,title={Before you move on}]
What does \hi{डुबोना} mean? (\textquotedblleft{}To dip; to immerse\textquotedblright{}.) Say it once more.
\end{tcolorbox}
\section[\texorpdfstring{\hi{बिखरना} (bikharnā)}{बिखरना (bikharnā)}]{\hi{बिखरना} (bikharnā) — to be scattered}

\label{lesson:HI-C333-bikharna}



\begin{quote}
Before the new one: say the Hindi for immersion (of an image), then the Hindi for to dip.
\end{quote}

\subsection*{You'll want to know: \hi{बिखरना}}
\textbf{\hi{बिखरना}} — \emph{bikharnā} — \textquotedblleft{}to be scattered\textquotedblright{}.

\begin{quote}
\hi{हवा} \hi{से} \hi{फूल} \hi{पूरे} \hi{घाट} \hi{पर} \hi{बिखर} \hi{गए।}
\end{quote}

\begin{quote}
\emph{havā se phūl pūre ghāṭ par bikhar gae.}
\end{quote}

\begin{quote}
The wind scattered the flowers all over the ghat.
\end{quote}

Past (it …): \hi{वह} \hi{माला} \hi{बिखर} \hi{गई}, \emph{vah mālā bikhar gaī}.

\begin{cousinweb}
\textbf{\hi{बिखरना}} (\emph{bikharnā}) is inherited from \emph{*viṣkirati}, \textquotedblleft{}scatters\textquotedblright{}, a form CDIAL reconstructs (CDIAL 11985). Flowers, beads or a crowd \textbf{\hi{बिखरते} \hi{हैं}}.
\end{cousinweb}

\begin{grammarlens}[title={What you've built}]
One more word for telling what happened.
\end{grammarlens}

\subsection*{Your turn}
\begin{itemize}
\raggedright
  \item \emph{Say it:} \emph{bikharnā}
  \item \emph{Say it:} \emph{bikharnā}, once more
  \item \emph{Say it:} \emph{havā se phūl pūre ghāṭ par bikhar gae.} — The wind scattered the flowers all over the ghat.
  \item \emph{Recall:} say the Hindi for immersion (of an image), then the Hindi for to dip, then say \emph{bikharnā} again
\end{itemize}

\begin{tcolorbox}[breakable,colback=teal!4,colframe=teal!35!black,title={Before you move on}]
What does \hi{बिखरना} mean? (\textquotedblleft{}To be scattered\textquotedblright{}.) Say it once more.
\end{tcolorbox}
\section[\texorpdfstring{\hi{दीपदान} (dīpdān)}{दीपदान (dīpdān)}]{\hi{दीपदान} (dīpdān) — the floating of lamps}

\label{lesson:HI-C333-dipdan}



\begin{quote}
Before the new one: say the Hindi for to dip, then the Hindi for to be scattered.
\end{quote}

\subsection*{You'll want to know: \hi{दीपदान}}
\textbf{\hi{दीपदान}} — \emph{dīpdān} — \textquotedblleft{}the floating of lamps\textquotedblright{}.

\begin{quote}
\hi{शाम} \hi{को} \hi{घाट} \hi{पर} \hi{दीपदान} \hi{हुआ।}
\end{quote}

\begin{quote}
\emph{śām ko ghāṭ par dīpdān huā.}
\end{quote}

\begin{quote}
In the evening lamps were floated on the river at the ghat.
\end{quote}

Masculine.

\begin{cousinweb}
\textbf{\hi{दीपदान}} (\emph{dīpdān}) joins Sanskrit \emph{dīpa}, \textquotedblleft{}a lamp\textquotedblright{}, and \emph{dāna}, \textquotedblleft{}a giving\textquotedblright{} (Monier-Williams): lamps given to the river, set afloat at dusk.
\end{cousinweb}

\begin{grammarlens}[title={What you've built}]
One more word for telling what happened.
\end{grammarlens}

\subsection*{Your turn}
\begin{itemize}
\raggedright
  \item \emph{Say it:} \emph{dīpdān}
  \item \emph{Say it:} \emph{dīpdān}, once more
  \item \emph{Say it:} \emph{śām ko ghāṭ par dīpdān huā.} — In the evening lamps were floated on the river at the ghat.
  \item \emph{Recall:} say the Hindi for to dip, then the Hindi for to be scattered, then say \emph{dīpdān} again
\end{itemize}

\begin{tcolorbox}[breakable,colback=teal!4,colframe=teal!35!black,title={Before you move on}]
What does \hi{दीपदान} mean? (\textquotedblleft{}The floating of lamps\textquotedblright{}.) Say it once more.
\end{tcolorbox}
